\section{Discussion}\label{sec:discussion}

\subsection{What the Scaling Step Measures}
Proposition~\ref{prop:identity} recasts per-epoch z-scoring as a measurement. Applied to a baseline-corrected epoch, it writes the post-stimulus mean amplitude into every sample, and most visibly into the pre-stimulus interval, where it is the whole of the signal's mean. The SHAPE observations follow the identity closely. A decoder fed $-m/\hat\sigma$ and a decoder fed the pre-stimulus mean of the z-scored epoch reach the same accuracy and agree on 95.9\% of their predictions. The time-resolved analysis shows the second consequence of the same step: every bin carries the epoch mean, so the decoding curve becomes flat in time and its late structure is attenuated. Under per-epoch scaling a decoding time course therefore reports the epoch mean at every latency and loses its meaning as a measurement of when information appears \cite{grootswagers2017decoding,king2014temporal}.

\subsection{Regimes of Leakage}
The averaging law organizes the observations across datasets. In single-trial decoding the leaked feature competes with trial noise, and its effect is small in most paradigms, from 0.4 to 3.7 points. In decoding of averaged data the trial noise shrinks as $1/k$, and the leak approaches the discriminability of the post-stimulus mean itself: 7 to 22 points across participants on ERP CORE and 20 points on SHAPE. Participant-level decoding of averaged ERPs is common in clinical and individual-differences work, where it is also the regime least protected by subject-grouped validation, because the leak lives inside each observation. The law also predicts where the leak is weak. When the condition difference is brief or changes sign within the epoch, as in the MMN, its contribution to the epoch mean is small, and the MMN shows no within-participant leak at any $k$. Across participants its z-scored statistic (0.553, $p=0.045$) sits close to its unscaled value (0.537, $p=0.070$), so its pre-stimulus interval is not a clean null region in the first place.

The prevalence audit places these regimes in current practice. Three of ten widely used codebases scale each sample by its own statistics in their default or reported pipelines, and 70\% of a fixed sample of recent decoding papers leave their input scaling unreported, which leaves the reader without the information needed to assess this leak.


\subsection{Re-Evaluation of the BIBM 2026 Reservoir Result}
The BIBM study characterized the ARSPI-Net temporal encoding stage as a perturbation-characterized measurement operator and reported BA 0.487 for its three-class affective decoding \cite{lane2026bibm}. The present analysis reproduces that value and identifies its source. The spike-count window of that study spans $-160$ to $+73$~ms, and its input epochs were z-scored over their full length. Under scaling that uses no post-stimulus sample, the same reservoir, readout, and protocol score 0.294 [0.262, 0.327], and the onset-truncated reservoir confirms that a causal system acquires the condition information before the stimulus arrives. The decodable content of the BIBM representation is therefore the leaked post-stimulus mean. The two propositions of that study, the chance floor of raw retention and the dependence of zero-fill scores on feature centering, are statements about scoring rules and hold for any input. The magnitudes it reported for them were measured on z-scored input and are re-measured under leak-free scaling as part of the program described below. The firing-rate and damage-spreading measurements describe the reservoir's dynamics under its input, while their link to decoding accuracy requires re-measurement with a post-stimulus read-out window and pre-stimulus scaling.

This re-evaluation was made possible by the transparency of the encoder. A fixed reservoir with an explicit read-out window can be truncated at onset, and that single manipulation separated genuine dynamics from leaked information. The same test applies to any temporal model, as the EEGNet analysis shows.

\subsection{A Reporting Rule}
Three practices follow directly from the theory. First, both the subtracted and the dividing statistic should come from data that contain no part of the response: the pre-stimulus interval of the same epoch, statistics fitted on training data and applied unchanged to test data, as EEG Conformer does \cite{song2023conformer}, or a causal running estimate \cite{schirrmeister2017deep}. Dropping the mean subtraction is not sufficient, because the scale route survives it (Table~\ref{tab:routes}). Second, every decoding pipeline should report its accuracy on the pre-stimulus interval alone, using the model itself or a second-order feature as the encoder, since a first-order check passes under scale-only normalization. Under a valid pipeline this accuracy sits inside its permutation null, so the check costs one additional run and detects this leak, the filtering displacements of \cite{vandriel2021highpass}, and any other transformation that moves post-stimulus information backward in time. Third, studies that decode averaged ERPs at the participant level should treat per-epoch scaling as a high-risk choice, because the averaging law places them in the regime where the leak is largest.

\subsection{Consequences for Online Systems}
Whole-epoch statistics are available only after the epoch has ended, so a decoder validated offline under per-epoch scaling has no causal implementation that reproduces its input. For a brain-computer interface or a closed-loop rehabilitation system, the practical consequence concerns timing claims more than headline accuracy. On the full epoch the scaling changes EEGNet accuracy by $+0.023$ [$-0.005$, $+0.052$], while on windows that end at or shortly after onset it accounts for the whole margin above chance. Early-decision and anticipatory decoders, and any estimate of the latency at which a class becomes separable, are therefore the analyses to re-examine, and the pre-stimulus null test is the criterion they should meet before deployment.

\subsection{Scope and Future Directions}
The present work characterizes trial-averaged affective ERPs in one cohort of 211 participants and five stimulus-locked paradigms of a public resource. The theory covers any time-domain location or scale statistic estimated from the epoch itself, and extending the measurements to spectral and connectivity features, to response-locked paradigms with a separately defined null interval, and to the window-level normalization used by EEG foundation models are natural next steps. The averaging law treats $\hat\sigma$ as constant, and a second-order account of the random denominator would sharpen the prediction at small $k$. For the ARSPI-Net program, the measurement operator of \cite{lane2026bibm} can now be re-measured with a post-stimulus read-out window and leak-free scaling, which gives its dynamical characterization a clean input on which to stand.
